%%%%%%%%%%  scan idx 287  |  folio 274  %%%%%%%%%%
\section*{\underline{EXPOSÉ IV}\\[3pt] \underline{\underline{GROUPES DE GROTHENDIECK DES TOPOS ANNELÉS}}}
% Exposé IV underlines parf (scan check, 10 Nov)
\renewcommand{\parf}{\underline{\mathrm{parf}}}
\begin{center}
\underline{par L. ILLUSIE}
\end{center}

\section*{1. Rappels et généralités sur les groupes de Grothendieck}

\textbf{1.1.} Soit \(C\) une catégorie \emph{triangulée}. Rappelons (I 6.3 et SGA 5 VIII 2) qu'une application \(f\) de \(\Ob C\) dans un groupe abélien \(G\) est dite \emph{additive} si l'on a \(f(E)=f(E')+f(E'')\) pour tout triangle distingué \(E'\longrightarrow E\longrightarrow E''\longrightarrow E'[1]\). Les applications additives de \(\Ob C\) dans \(G\) forment un groupe abélien, qui dépend fonctoriellement de \(G\). Le foncteur ainsi obtenu est représenté par un groupe abélien \(k(C)\) et une application additive universelle \(\cl_C:\Ob C\longrightarrow k(C)\) (notée \(\cl\) quand il n'y a pas de confusion à craindre).

Le groupe \(k(C)\) dépend fonctoriellement de \(C\) vis-à-vis des foncteurs exacts. Si \(C,C'\) sont des catégories triangulées, deux foncteurs exacts isomorphes de \(C\) dans \(C'\) induisent le même homomorphisme \(k(C)\to k(C')\); en particulier, si \(u:C\to C'\) est un foncteur exact qui est une équivalence de catégories, \(k(u):k(C)\to k(C')\) est un isomorphisme.

\textbf{Lemme 1.2.} Soient \(C\) une catégorie triangulée, et \(L\in\Ob C\).
\begin{enumerate}[label=\alph*)]
\item Pour tout \(n\in\mathbb Z\), on a
\[
 \cl(L[n])=(-1)^n\cl(L).
\]
Si \(L',L''\) sont des objets de \(C\) tels qu'on ait \(L\simeq L'\oplus L''\), alors on a \(\cl(L)=\cl(L')+\cl(L'')\).
\item Supposons qu'il existe \(n\in\mathbb Z\) tel que \(\coprod_{i\geq0} L[2ni]\) soit représentable. Alors on a \(\cl(L)=0\).
\end{enumerate}

\emph{Preuve.} L'assertion a) découle directement des définitions (cf. SGA 5 VIII 2). Pour b), il suffit de
remarquer qu'on a
